\sectionguard
\section*{References}

{\TriptychReferenceFont
\clubpenalty=10000
\widowpenalty=10000
\raggedright
\begin{itemize}[leftmargin=1.1em,labelsep=0.35em,itemsep=0.10em,topsep=0.15em,parsep=0pt]
\item \emph{Missale Romanum}, editio typica (Vatican: Typis Polyglottis
Vaticanis, 1962), \latin{Dominica Decima Quinta post Pentecosten},
pp.~396--397, nos.~1582--1591,
\sourceurl{https://media.churchmusicassociation.org/pdf/missale62.pdf}{CMAA facsimile};
Benziger, New York, 1962, editio iuxta typicam, pp.~390--392,
\sourceurl{https://archive.org/details/MissaleRomanum1962RomanMissalColorLatin}{independent edition};
Tours: Alfred Mame, 1922, editio quarta iuxta typicam Vaticanam, the
same formulary, pp.~446--448. The Mame edition supplies the recorded
United States public-domain antecedent; the Vatican edition controls the
appointed text.

\item \emph{The Roman Missal translated into the English language for the
use of the laity} (Philadelphia: Eugene Cummiskey, 1861), Fifteenth Sunday
after Pentecost, Collect, Secret and Postcommunion. The continuation's XVI
running head is erroneous; the three incipits and formulary sequence
identify the correspondence. Historical English witness, not an approved
1962 translation.

\item \emph{The Bible, Douay-Rheims, Complete}, Challoner revision,
\sourceurl{https://www.gutenberg.org/ebooks/1581}{Project Gutenberg 1581}:
Pss.~39,46,85,91,94; Gal.~1--6; Luke~7, with 1:68--79;
John~6; Exod.~4:29--31 and 34:5--9; 3~Kings~17:17--24;
4~Kings~4:8--37; Heb.~3--4 and 10:5--10.
The cultural gallery's quoted comparison is the separate
\emph{Douay--Rheims, American Edition of 1899}, eBible.org engDRA,
Ps.~39:4, Gal.~6:2,7--8 and Luke~7:12,14--15. Both historical English
witnesses are public domain in the United States.

\item Clementine Vulgate, eBible.org latVUC electronic witness,
host-identified as the 1598 text with Glossa Ordinaria in the Migne 1880
edition, \sourceurl{https://eBible.org/Scriptures/latVUC_usfm.zip}{USFM distribution}:
Pss.~39,46,85,91,94; Gal.~5:13--6:10; Luke~7:11--17;
John~6:49--60; 3~Kings~17:17--24; Heb.~10:5--10.
Biblical text, not the accompanying Glossa, supplies the comparisons.

\item Augustine, \emph{Enarrationes in Psalmos}, 39.2--5, 46.2--3,
85.1--6, 91.1--4 and 94.5--6,10; J.~E. Tweed's English,
NPNF first series~8 (Buffalo, 1888), New Advent revised presentation:
\sourceurl{https://www.newadvent.org/fathers/1801040.htm}{Psalm 40},
\sourceurl{https://www.newadvent.org/fathers/1801047.htm}{47},
\sourceurl{https://www.newadvent.org/fathers/1801086.htm}{86},
\sourceurl{https://www.newadvent.org/fathers/1801092.htm}{92},
\sourceurl{https://www.newadvent.org/fathers/1801095.htm}{95}.
The online headings use modern psalm numbers; the citations above use
Vulgate numbers.

\item Augustine, \emph{Sermo}~98, §§1--7, especially 4--5;
R.~G. MacMullen's English, NPNF first series~6 (Buffalo, 1888),
\sourceurl{https://www.newadvent.org/fathers/160348.htm}{New Testament Sermon 48}.
The web heading's Luke~7:2 is erroneous; §4 narrates Naim.
\emph{In Iohannis Evangelium tractatus} 26.10--20, especially
13,15--18; John Gibb's English, NPNF first series~7 (Buffalo, 1888),
\sourceurl{https://www.newadvent.org/fathers/1701026.htm}{Tractate 26}.

\item Augustine, \emph{Sermo Frangipane V}, §§2--3 and §4 opening,
in Germain Morin, ed., \emph{Sancti Augustini Sermones post Maurinos
reperti}, \emph{Miscellanea Agostiniana} I (Rome: Typis Polyglottis
Vaticanis, 1930), pp.~213--215; heading and editorial introduction,
p.~212, \sourceurl{https://archive.org/download/bwb_W9-DJJ-981_1/bwb_W9-DJJ-981_1.pdf}{complete facsimile}.
The Latin locus supplies the argument also circulated under Sermo~163B;
Morin's p.~214 apparatus marks a repeated charity formula as a possible
gloss and his sermon year remains tentative.

\item Ambrose, \emph{Expositio Evangelii secundum Lucam} V.89--92,
Migne 1845 Latin edition, continuous Corpus Corporum-derived
\sourceurl{https://la.wikisource.org/w/index.php?title=Expositio_Evangelii_secundum_Lucam_(Ambrosius)/5\&action=raw}{Wikisource book V transcription},
column marker 1377 within §91. Bede, \emph{In Lucae Evangelium expositio}
II.7, on Luke~7:11--16, PL~92:417B--419A, especially 418A--D,
\sourceurl{https://la.wikisource.org/w/index.php?title=In_Evangelium_S._Lucae_(Beda)\&action=raw}{continuous Latin transcription}.
Both are attributed transcriptions, not new facsimile or critical-edition
collations.

\item Cyril of Alexandria, \emph{A Commentary upon the Gospel according
to S. Luke}, trans. R.~Payne Smith, I (Oxford, 1859), Sermon~XXXVI,
pp.~132--135, \sourceurl{https://upload.wikimedia.org/wikipedia/commons/9/90/Commentary_upon_the_Gospel_of_Luke_by_st._Cyril_of_Alexandria._Volume_1.pdf}{volume I facsimile}.
The editor supplies Greek-catena text for a missing Syriac folio at
p.~132 and further Greek fragments in p.~135's note.
\emph{Commentary on John} IV.2, trans. P.~E. Pusey (Oxford, 1874),
Library of Fathers 43; pp.~409--414 on the appointed clause and
418--419 on the following verse's Naim comparison,
\sourceurl{https://www.tertullian.org/fathers/cyril_on_john_04_book4.htm}{continuous English book IV}.
This is Pusey's portion of the 1874/1885 collection, not Payne Smith's
Luke translation or Randell's later completion.

\item John Chrysostom, \emph{Commentary on Galatians}, chapters~5--6,
lemmata 5:25--6:10, trans. Gross Alexander, NPNF first series~13
(Buffalo, 1889), \sourceurl{https://www.newadvent.org/fathers/23105.htm}{chapter 5}
and \sourceurl{https://www.newadvent.org/fathers/23106.htm}{chapter 6};
\emph{Homilies on John} 46.1--4, trans. Charles Marriott, NPNF first
series~14 (Buffalo, 1889),
\sourceurl{https://www.newadvent.org/fathers/240146.htm}{Homily 46}.
Historical English witnesses; modern host presentation is a separate
rights object.

\item Theodoret of Cyrus, \emph{Interpretatio in Psalmos}, PG~80
(Paris, 1860): Ps.~39, cols.~1151B--1154C; Ps.~46,
1205D--1208B; Ps.~85,
1553C--1556B; Ps.~91, 1615C--1618A; Ps.~94,
1639B--1642A, exact all-earth lemma at 1641A and 1642A,
\sourceurl{https://upload.wikimedia.org/wikipedia/commons/f/fd/Patrologia_Graeca_Vol._080.pdf}{Greek--Latin facsimile}.
The Ps.~39 apparatus at 1154 reports a commentary-order difficulty.
\emph{Interpretatio Epistolae ad Galatas} 6:1--10, PG~82
(Paris, 1864), cols.~499A--502B,
\sourceurl{https://upload.wikimedia.org/wikipedia/commons/f/f2/Patrologia_Graeca_Vol._082.pdf}{Greek--Latin facsimile}.

\item Robert Bellarmine, \emph{A Commentary on the Book of Psalms},
trans. John O'Sullivan (Dublin and London: James Duffy, 1866):
Ps.~39, pp.~118--119; Ps.~46, 143--145, especially p.~144,
edition v.~2; Ps.~85, 259--260; Ps.~91,
289--290; Ps.~94, 298--299;
\sourceurl{https://upload.wikimedia.org/wikipedia/commons/9/9d/Commentaryonbook0000bell.pdf}{facsimile}.
The translator's p.~v explains omitted philology and replacement of
Bellarmine's psalm prefaces by Douay headings.

\item Thomas Aquinas, \emph{Super Epistolam B. Pauli ad Galatas lectura}
6.1--2, Turin 1953 Latin base, Busa transcription revised by Enrique
Alarcón, \sourceurl{https://www.corpusthomisticum.org/cgl.html}{Corpus Thomisticum},
electronic units [87781--87782];
\emph{Super Evangelium S. Ioannis lectura} 6.6, Turin 1952 base,
\sourceurl{https://www.corpusthomisticum.org/cih06.html}{electronic unit [87463]}.
The bracketed identifiers are electronic text units, not printed Marietti
paragraphs; Latin arguments are paraphrased, with no Larcher English
quotation. \emph{Summa theologiae} III.79.1--3,6, Fathers of the
English Dominican Province, second revised edition (1920),
\sourceurl{https://www.newadvent.org/summa/4079.htm}{complete question}.
\emph{In psalmos Davidis expositio},
\emph{Super Psalmo}~46, nn.~1--4, especially n.~2 [87252],
\emph{Reportatio Reginaldi de Piperno}, Parma 1863 base, Busa
transcription revised by Enrique Alarcón,
\sourceurl{https://www.corpusthomisticum.org/cps41.html\#87252}{Corpus Thomisticum}.

\item H.~A. Wilson, ed., \emph{The Gelasian Sacramentary}
(Oxford: Clarendon, 1894), III.xi, \emph{Item alia missa}, p.~230,
\sourceurl{https://archive.org/download/gelasiansacrame00wilsgoog/gelasiansacrame00wilsgoog.pdf}{facsimile};
\emph{The Gregorian Sacramentary under Charles the Great}, Henry Bradshaw
Society~49 (London, 1915), supplement XXXII--XXXIIII,
pp.~173--174; introduction pp.~xvii--xviii, xxi and xliv--xlv,
\sourceurl{https://archive.org/download/gregoriansacrame00cath/gregoriansacrame00cath.pdf}{facsimile}.

\item \emph{The Liturgical Year}, continuation published under Abbot
Guéranger's name, trans. Laurence Shepherd, \emph{Time after Pentecost}
II (overall XI), second edition (Stanbrook Abbey, Worcester;
London: Burns \& Oates, R.~\& T. Washbourne and Art Book Company;
United States: Benziger, 1909), pp.~348,350--351,354--355,
\sourceurl{https://archive.org/download/V11TheLiturgicalYear/V11TheLiturgicalYear.pdf}{facsimile}.
The title page and continuation preface control this imprint rather than
the catalog's 1900/Duffy label; the series attribution does not establish
Guéranger's personal authorship of the continuation.

\item Ildefonso Schuster, \emph{The Sacramentary (Liber Sacramentorum)},
trans. Arthur Levelis-Marke, III (London: Burns Oates \& Washbourne,
1927), Fifteenth Sunday, pp.~139--141; following heading, p.~142.

\item Scripture chronology witnesses: USCCB, \emph{New American Bible,
Revised Edition}, \sourceurl{https://bible.usccb.org/bible/psalms/0}{Psalms introduction};
\emph{Catholic Encyclopedia}, ``Epistle to the Galatians,'' VI (1909),
``St. Paul,'' XI (1911), ``The New Testament,'' XIV (1912),
``Gospel of St. John,'' VIII (1910), composition notices;
A.~J. Maas, ``Jesus Christ,'' VIII (1910), III.B.3, Third and Fifth
Journeys, pp.~378,380. Cornelius a Lapide, \emph{Commentaria in omnes
divi Pauli epistolas} (Antwerp, 1614), \emph{Prooemium} and
Galatians \emph{Argumentum}. Pontifical Biblical Commission,
\emph{De auctore, de tempore compositionis et de historica veritate
Evangeliorum secundum Marcum et secundum Lucam}, responsa I--IX,
\emph{AAS}~4 (1912), pp.~463--465.

\item Oscar Wilde, \emph{The Importance of Being Earnest: A Trivial
Comedy for Serious People}, three-act text, Act~II, Miss Prism's two
sow/reap speeches and Chasuble's appeal for charity;
\sourceurl{https://www.gutenberg.org/ebooks/844}{Gutenberg 844}, David Price
transcription, updated 10~November 2025. The header identifies no print
copy-text; act and dialogue locate the passages.

\item William Lloyd Garrison, ``Guilt of the North---The National Compact,''
in Richard Sutton Rust, comp., \emph{Freedom's Gift: or, Sentiments of
the Free} (Hartford: S.~S. Cowles, 1840), pp.~34--40,
especially p.~38, paragraph beginning ``Divided, we fall,''
\sourceurl{https://tile.loc.gov/storage-services/public/gdcmassbookdig/freedomsgiftorse00rust/freedomsgiftorse00rust.pdf}{Library of Congress facsimile},
PDF leaf~46; \sourceurl{https://www.loc.gov/item/30020058/}{item 30020058},
identity and Rights \& Access.

\item U2, ``40,'' track~10 on \emph{War} (1983),
\sourceurl{https://www.u2.com/pages/song/40}{official song page} and
\sourceurl{https://www.u2.com/products/war}{album page}; U2.com,
\sourceurl{https://www.u2.com/blogs/news/this-is-40}{``This is 40''},
5~March 2024, report of 2~March's final U2:UV show, especially the
paragraph beginning ``As the band left the stage one by one.'' Lyrics
and site prose remain protected; the gallery quotes eight lyric words.

\item Lothar Warneke, dir., \emph{Einer trage des anderen Last\ldots}
(DEFA-Studio für Spielfilme, production 1987; premiere 28~January 1988).
Philip Zengel, \sourceurl{https://www.defa-stiftung.de/stiftung/aktuelles/film-des-monats/einer-trage-des-anderen-last/}{DEFA Film of the Month},
October 2021: \emph{Kurzinhalt}, \emph{Produktionsnotizen},
\emph{Gleichgewichtigkeit der Filmfiguren} and following parable discussion;
UMass DEFA Film Library, \sourceurl{https://www.umass.edu/defa/film/3523/}{catalog 3523},
Synopsis and Commentary. ERF Bibleserver,
\sourceurl{https://www.bibleserver.com/LUT/Galater6}{Lutherbibel 2017, Galatians 6:2},
©~2016 Deutsche Bibelgesellschaft; later wording comparison only.

\item Charles Dickens, \emph{The Seven Poor Travellers}, II,
``The Story of Richard Doubledick,'' visit to Taunton's mother at Frome
and subsequent recovery; David Price's transcription from
\emph{Christmas Stories} (Chapman and Hall, 1894),
\sourceurl{https://www.gutenberg.org/ebooks/1392}{Gutenberg 1392}, updated
31~December 2020. The identifying paragraphs begin ``Though he was weak
and suffered pain,'' ``It was a Sunday evening,'' and ``But that night.''
The new mother--son relation continues in the later recovery passage.
\end{itemize}

\noindent Web witnesses were accessed 5~September 2026; appointed-text and
chronology checks retain their separate recorded dates.
}
